\chapter*{Nota de método}
\addcontentsline{toc}{chapter}{Nota de método}

Este documento pertenece al género de las \emph{Implicaciones del Paquete Económico}: la
arquitectura en tres partes sobre la restricción presupuestal, la plantilla de capítulo, la
densidad de cuadros. Se aparta de él en seis lugares, y los seis se declaran aquí.

\section*{Uno. Está escrito bajo régimen de simulación en vivo}

No se usó ninguna pieza aprobada del ejercicio 2026. Ni la Ley de Ingresos aprobada, ni el
Presupuesto de Egresos aprobado, ni sus analíticos, ni la Cuenta Pública de 2025, ni informes
trimestrales posteriores a la entrega del paquete. El documento del género sobre el mismo
paquete permaneció en cuarentena, sin abrirse, durante toda la redacción.

\textbf{Por qué.} Un documento escrito con el presupuesto ya aprobado es más exacto y es
irrepetible el día que llega el paquete. Este ejercicio existe para saber qué se puede
afirmar con lo que hay en la mano esa mañana.

\section*{Dos. Dos líneas de comparación, siempre etiquetadas}

\begin{list}{}{\setlength{\leftmargin}{1.2em}}
\item[] \textbf{Línea P, primaria.} Proyecto 2026 contra \emph{proyecto} 2025. Ex ante contra
ex ante. Es la comparación disponible el día de la entrega y la que nadie más construye.
\item[] \textbf{Línea G, del género.} Proyecto 2026 contra \emph{aprobado} 2025. Es la que usa
la discusión pública, y mezcla un objeto ex ante con uno ex post.
\end{list}

Toda columna comparativa lleva su letra. \textbf{Nunca conviven sin etiqueta.} Donde difieren de
forma material, la diferencia mide lo que la Cámara cambió el año anterior, y se reporta como
hallazgo.

Hay una razón concreta para esta manía. La misma función educativa de 2026, con el mismo
perímetro y el mismo deflactor, crece \textbf{3,4\,\%} en la línea P y \textbf{2,2\,\%} en la
línea G. Las dos son verdaderas y el punto y dos décimas de diferencia no es ruido: es lo que
la Cámara añadió a la función educativa en diciembre de 2024 ---12\,641,1 millones de pesos
entre el proyecto y el aprobado de 2025--- reapareciendo como una base más alta.

El precedente que obliga a etiquetar es del \emph{ejercicio anterior}, y ahí la elección de
línea cambiaba el signo: la función educativa del proyecto de 2025 crecía \textbf{0,81\,\%}
contra el proyecto de 2024, \textbf{0,70\,\%} de aprobado contra aprobado, y \textbf{caía
0,47\,\%} en la línea G. Ninguna de las tres era falsa.

\section*{Tres. Toda diferencia en millones de pesos declara si es nominal o real}

Los cuadros llevan las dos columnas. El deflactor es el del PIB que declara el CGPE 2026:
\textbf{4,8\,\%}. Nunca se deriva un deflactor de agregados publicados redondeados.

\begin{recuadro}{Por qué el deflactor del PIB, y cuándo no}
\textbf{La elección de índice no es una preferencia de estilo: es parte de la definición del
objeto que se mide.} Este documento usa dos, y los usa en dominios que no se tocan.

\textbf{Deflactor implícito del PIB} ---4,8\,\% para 2026, publicado por el CGPE--- para
\emph{todo agregado fiscal y toda razón a PIB}. La razón es de consistencia: el denominador de
esas razones es el PIB nominal, y deflactar el numerador con un índice distinto del que
gobierna el denominador produce una serie que no es de nada. Es también el índice con el que
la Secretaría construye el límite de gasto corriente estructural, de modo que la regla fiscal
y este documento miden con la misma vara.

\textbf{Índice de precios al consumidor} ---y sólo él--- cuando la pregunta es el
\emph{poder de compra de una prestación individual}: cuánto compra la pensión de una persona,
cuánto vale una beca en el bolsillo de quien la recibe. Ahí la pregunta es de bienestar y el
consumidor es el sujeto, así que el índice del consumidor es el correcto. \textbf{Ese uso no
aparece en ningún cuadro presupuestal de este documento}, para que nadie sume dos columnas
construidas con varas distintas.

\textbf{Sensibilidad, una sola línea y aquí.} Con la inflación promedio esperada para 2026 en
lugar del deflactor del PIB, toda variación real de este documento sería aproximadamente
\textbf{un punto porcentual mayor}. Esa es exactamente la distancia que separa nuestras
cifras de las del documento del género, que usa 1,0365 sin declararlo ---tercer año
consecutivo---. \textbf{La línea vive en este recuadro y no se repite por capítulo}, para que
ningún cuadro mezcle las dos convenciones.

\textbf{Lo prohibido, que este proyecto ha visto tres veces en el género:} deflactar con un
índice y rotular la columna con el nombre del otro; presentar una columna del ejercicio
anterior ya deflactada bajo el rótulo del año anterior; y derivar un deflactor de dos
agregados publicados redondeados.
\end{recuadro}

\section*{Cuatro. Hay una capa demográfica declarada}

Los denominadores de población no están en el paquete y sí están al alcance. Se usan, en
cuadros \emph{separados} de los presupuestales, con el tier \texttt{externa\_demografica} y
con fuente, año de referencia y liga en el archivo de fuentes. Ninguna cifra del paquete
aparece junto a un denominador externo en la misma línea sin marca.

Donde un denominador no se consiguió, \textbf{se declara y no se calcula}. Este documento no
publica gasto por afiliado a IMSS o ISSSTE, ni gasto por alumno matriculado, y dice por qué.

\section*{Cinco. El capítulo de deuda descompone, y declara lo que no puede identificar}

El cambio de la razón del saldo histórico a PIB se reparte en cuatro variables: crecimiento,
inflación por el denominador, tipo de cambio y el requerimiento del año. Las cuatro suman el
cambio observado. La \emph{compensación por inflación} que el sector público paga sobre su
deuda no es un término del marco sino una lectura del requerimiento, y \textbf{el paquete no
identifica qué índice de precios le corresponde}. Se presenta como rango con su cota, nunca
como punto.

\section*{Seis. Hay un registro de lo que no se pudo verificar}

Al final del documento, y no solo en la bitácora de trabajo. Incluye las verificaciones que
son imposibles bajo régimen en vivo por construcción.

\section*{Convenciones}

\begin{list}{}{\setlength{\leftmargin}{1.2em}\setlength{\itemsep}{1pt}}
\item[] \textbf{Nomenclatura.} RFSP y SHRFSP; las grafías RFSPF y SHRFSPF designan el mismo
objeto y la equivalencia se declara aquí una sola vez.
\item[] \textbf{Perímetros.} Declarados antes del primer cuadro de cada capítulo, enumerables
por clave, y constantes dentro del capítulo.
\item[] \textbf{Claves, no nombres.} Toda agregación por ramo, programa o unidad responsable
se hace por su clave numérica. Y cuando la clave misma cambia de prefijo, por su número: la
Pensión Mujeres Bienestar es U316 en 2025 y S316 en 2026.
\item[] \textbf{Ramo y función son objetos distintos.} Una caída de ramo no es una caída de la
función que su nombre sugiere.
\item[] \textbf{Umbrales.} 0,5\,\% en niveles y 0,1 puntos en tasas y razones.
\item[] \textbf{Definiciones de balance y deuda} adjudicadas contra \emph{Balance Fiscal en
México} (SHCP, abril de 2023), nunca por criterio propio.
\item[] \textbf{Tier de cada cifra:} \texttt{oficial\_primaria}, \texttt{derivada\_ciep},
\texttt{autoral\_ited}, \texttt{externa\_demografica}. Registrado en el archivo de fuentes.
\end{list}

\section*{Validación}

Treinta y tres identidades contables y de coherencia entre anexos, todas cerradas, corridas
después de cada cambio en los datos. El bruto del proyecto ---11\,746\,796,8 millones de
pesos--- se obtuvo por tres rutas independientes en dos servidores distintos y coincide al
peso con el Anexo 1 del decreto. \textbf{Se valida por identidad, no por relectura:} un
dígito mal leído rompe una suma, y una relectura lo repite.
\clearpage
